\section{A reciprocal fractional cube root and its generating function}

Kimberling's sequence A184534 is defined, for $n\ge1$, by
\[
 a(n)=\left\lfloor\frac1{\{(n^3+4)^{1/3}\}}\right\rfloor,
\]
where braces denote fractional part. Barker recorded an empirical rational
generating function and a separate empirical closed form in 2012
\cite{oeis184534}. We prove the generating function and identify the
one-index displacement in the printed closed form.

\begin{theorem}
The source sequence satisfies
\[
 a(1)=1,\qquad a(2)=3,\qquad a(3)=7,\qquad
 a(n)=\left\lfloor\frac{3n^2}{4}\right\rfloor\quad(n\ge4).
\]
Its ordinary generating function is exactly
\begin{equation}\label{a184534:gf}
 \sum_{n\ge1}a(n)x^n=
 \frac{x(x^6-2x^5+x^4-x^2-x-1)}{(x-1)^3(x+1)}.
\end{equation}
For every $n\ge8$,
\[
 a(n)=2a(n-1)-2a(n-3)+a(n-4).
\]
\end{theorem}

\begin{lemma}\label{a184534:general}
For positive integers $h\le n$,
\[
 \left\lfloor\frac1{\{(n^3+h)^{1/3}\}}\right\rfloor
 =\left\lfloor\frac{3n^2}{h}\right\rfloor.
\]
\end{lemma}
\begin{proof}
Put $q=\lfloor3n^2/h\rfloor$. Then $q\ge3n>0$,
$hq\le3n^2$, and, by integrality, $h(q+1)\ge3n^2+1$.
First,
\[
 (nq+1)^3-(n^3+h)q^3
 =q^2(3n^2-hq)+3nq+1>0.
\]
Second,
\begin{align*}
 &(n^3+h)(q+1)^3-(n(q+1)+1)^3\\
 &\quad=(q+1)^2\bigl(h(q+1)-3n^2\bigr)-3n(q+1)-1\\
 &\quad\ge(q-3n)(q+1)+q>0.
\end{align*}
Strict monotonicity of cubing therefore gives
\[
 n+\frac1{q+1}<(n^3+h)^{1/3}<n+\frac1q.
\]
In particular the root has integer part $n$. Its positive fractional
part lies strictly between $1/(q+1)$ and $1/q$, so its reciprocal lies
strictly between $q$ and $q+1$. The floor is $q$.
\end{proof}

\begin{proof}[Proof of the theorem]
Apply Lemma~\ref{a184534:general} with $h=4$ for every $n\ge4$.
The three earlier values follow from the same two cubic comparisons,
using the following exact integer certificates:
\[
\begin{array}{c|c|c|c}
n&q&(n^3+4)q^3<(nq+1)^3&
(n(q+1)+1)^3<(n^3+4)(q+1)^3\\\hline
1&1&5<8&27<40\\
2&3&324<343&729<768\\
3&7&10633<10648&15625<15872
\end{array}
\]
These comparisons also show the root has integer part $n$ in each
case; no numerical approximation to a root is required.

For every natural $n$, the identity
$\lfloor3n^2/4\rfloor=3\lfloor n^2/4\rfloor$ follows from
$n^2\equiv0$ or $1\pmod4$. The only positive-index corrections to this
formula are one additional unit at $n=1$ and at $n=3$. Splitting into
even and odd indices gives
\[
 \sum_{n\ge0}\left\lfloor\frac{n^2}4\right\rfloor x^n
 =\frac{x^2}{(1-x)^3(1+x)}.
\]
Consequently
\[
 \sum_{n\ge1}a(n)x^n
 =x+x^3+\frac{3x^2}{(1-x)^3(1+x)}.
\]
Combining the terms gives~\eqref{a184534:gf}. The normalized denominator
is $1-2x+2x^3-x^4$ and its numerator has degree seven. Coefficient
comparison proves the recurrence from newest index eight onward.
\end{proof}

\paragraph{The source's closed-form indexing.}
The expression printed as an empirical value of $a(n)$ is
\[
 \frac{3(1-(-1)^n+4n+2n^2)}8\qquad(n>2).
\]
It equals $12$ at $n=3$, while $a(3)=7$. The expression is instead
$a(n+1)$ for every $n\ge3$. The corrected same-index expression is
\[
 a(n)=\frac{3(2n^2-1+(-1)^n)}8\qquad(n\ge4).
\]
The recurrence cutoff also matters: at newest index seven its right
side is $37$, whereas $a(7)=36$.

\paragraph{Formal conclusion and attribution.}
Four Lean modules prove the literal real-root and fractional-part
statement, the initial cubic certificates, the all-index sequence
formula, the recurrence, and the exact formal generating function in
$\mathbb Z[[X]]$. The final theorem \texttt{A184534.full\_resolution}
also identifies every positive source coefficient and the zero
constant coefficient. The source's denominator has constant term
$-1$, so the formal quotient is unambiguous. A separate Lean theorem
certifies the printed closed-form discrepancy at $n=3$.

The eventual expression already occurs in sequence generators, and
the root-comparison method is elementary. Those expressions and
methods are not claimed as new discoveries. The contribution counted
here is one complete proof of the recorded empirical generating-function
assertion, with the indexing issue documented. The general array lemma
and the recurrence are parts of the same package. The separate bounded
prior-work record found no earlier proof connecting the literal source
to this identity and does not establish worldwide priority.
